\documentclass[10pt,a4paper]{amsart}
\usepackage[margin=19mm]{geometry}
\usepackage{amsmath,amssymb,mathtools}
\usepackage[scr=rsfs,cal=euler]{mathalfa}
\usepackage{xfrac}
\usepackage[unicode,hidelinks,bookmarks=false]{hyperref}
\newcommand{\ie}{i.e.\ }
\newcommand{\cf}{cf.\ }
\newcommand{\resp}{respectively\ }
\newcommand{\Subsection}{\subsection}
\providecommand{\nobreakdash}{\nobreak\hbox{-}}
\setlength{\emergencystretch}{2em}
\multlinegap0pt
\usepackage{mathtools}
\usepackage{tikz}
\usepackage{amssymb}
\usepackage{xcolor}
\newtheorem{theorem}{Theorem}
\newtheorem{lemma}{Lemma}
\newtheorem{proposition}{Proposition}
\def\mc#1{\mathcal{#1}}
\title{Stability of Wasserstein projections in convex order via metric extrapolation}
\author{Jakwang Kim, Young-Heon Kim, Andrea Natale}
\date{Source: arXiv:2507.19221v3 (CC BY 4.0)}
\begin{document}
\maketitle

\noindent\emph{Source and licence.} Stability of Wasserstein projections in convex order via metric extrapolation. Version arXiv:2507.19221v3 (CC BY 4.0), distributed by arXiv under the Creative Commons Attribution 4.0 International licence, http://creativecommons.org/licenses/by/4.0/, as recorded in the arXiv licence metadata for this exact version and shown on the abstract page https://arxiv.org/abs/2507.19221v3. Full licence notice: \texttt{licenses/arxiv-2507.19221-CC-BY-4.0.txt}.\par\smallskip

\noindent\textit{Editorial scope.} Counted unit: the article's proposition prop:extralip (source0.tex lines 269--273). The article's complete original proof of it is delivered with it (source0.tex lines 274--293, one \texttt{proof} environment of 20 lines). No mathematics is written, repaired or re-derived: the statement and the proof are the article's own lines, in the article's own notation, and no retained line is substituted or re-flowed.  Retained uncounted as the source-local context the proof needs: lines 142--158; lines 218--220; lines 227--233; lines 251--253.  Nothing was omitted from any retained span, so the package carries no omission record.  No retained line contains a citation, so the wrapper carries no bibliography and no bibliography entry.  No retained line contains a figure environment, an image-inclusion command or drawing code, so no image is delivered and the package declares no asset dependency.  Editorial formatting only, in addition: an AMS \texttt{amsart} preamble that restates the article's own declarations and macros used by the retained lines, namely the environments \texttt{theorem}, \texttt{lemma}, \texttt{proposition} on separate counters, together with the standard packages those lines need.  The retained environments print in this document as Theorem 1, Lemma 1, Proposition 1 in order of appearance, whereas the article prints the same items with the numbering of its own full document.  The complete native source of the article as distributed in the arXiv e-print for this exact version is bundled unchanged as \texttt{sources/182197/source0.tex}.
\begin{theorem}\label{th:projextra} Let $t>1$, 
$\theta \coloneqq (t-1)/t$, and let $\nu_0,\nu_1\in \mc{P}_2(\mathbb{R}^d)$  be arbitrary measures. The following holds:
\begin{enumerate}
    \item There exists a unique minimizing geodesic $\gamma:[0,t] \rightarrow \mc{P}_2(\mathbb{R}^d)$ such that \[
    \gamma(1) = \nu_1~\text{ and }~\gamma(t) = \mathscr{E}^t(\nu_0,\nu_1) .\]
    Moreover $\gamma(0) = \mathscr{P}^{cx}_{\preceq \nu_0}(D^{1/\theta}_\# \nu_1)$  and there exists a lower semi-continuous $\theta$-strongly convex potential $u:\mathbb{R}^d \rightarrow \mathbb{R} \cup \{+\infty\}$ such that $\gamma(s) = (s\mathrm{Id}+ (1-s)\nabla u^*)_\#\nu_1$ for all $s\in[0,t]$. 
    { Here $u^*$ is the Legendre transform of $u$.}

  

\item If  $\nu_0$ is absolutely continuous, there exists a unique minimizing geodesic $\tilde{\gamma}:[0,t]\rightarrow \mc{P}_2(\mathbb{R}^d)$ such that
\[
\tilde{\gamma}(0)=\nu_0 \quad \text{and}\quad \tilde{\gamma}(t) = \mathscr{E}^t(\nu_0,\nu_1)\,.
\]
Moreover,  $\tilde{\gamma}(1) = \mathscr{P}^{cx}_{\nu_1 \preceq}(D^\theta_\#\nu_0)$ and $\tilde{\gamma}(s) = ((1-s)\mathrm{Id}+ s\nabla u)_\#\nu_0$ for all $s\in[0,t]$, with the same potential $u$ as in point $(1)$.    
\end{enumerate}
\end{theorem}

\begin{equation}\label{eq:projl2}
\inf_{X\in L^2(\rho_0)} \left\{ \frac{1}{2}\|X - X_\mu \|^2_{ L^2(\rho_0)} \; :\; X_\# \rho_0  \preceq \nu \right\}\,.
\end{equation}

\begin{lemma}\label{lem:lift} 
{For  $\rho_0 \in \mc{P}_2^{ac}(\mathbb{R}^d)$ and $\mu, \nu  \in \mc{P}_2(\mathbb{R}^d)$,}
Problem \eqref{eq:projl2} admits a unique solution $X^*\in L^2(\rho_0)$. Furthermore, $X^* = T\circ X_\mu$ where $T\in L^2(\mu)$ is the unique optimal transport map from $\mu$ to $\mu^* \coloneqq \mathscr{P}_{\preceq \nu}^{cx}(\mu)$. In particular
\begin{equation}\label{eq:equality}
X^*_\# \rho_0 =  \mu^* \quad \text{and} \quad W_2(\mu,\mu^*) = \|X_\mu - X^*\|_{L^2(\rho_0)}\,. 
\end{equation}
\end{lemma}

\begin{equation}\label{eq:stabproj}
W_2(\mathscr{P}_{\preceq \nu}^{cx} (\mu),\mathscr{P}_{\preceq \nu}^{cx} (\tilde{\mu}))\leq \|X^* - \tilde{X}^*\|_{L^2(\rho_0)} \leq \| X_\mu - X_{\tilde{\mu}}\|_{L^2(\rho_0)}.
\end{equation}

\begin{proposition}\label{prop:extralip} Given  $\nu_0 \in \mc{P}_2(\mathbb{R}^d)$, for any $\nu_1, \tilde{\nu}_1\in\mc{P}_2(\mathbb{R}^d)$ and $t>1$,
\begin{equation}\label{eq:extralip}
W_2(\mathscr{E}^t(\nu_0,\nu_1),\mathscr{E}^t(\nu_0,\tilde{\nu}_1)) \leq 2t W_2(\nu_1,\tilde{\nu}_1).
\end{equation}
\end{proposition}

\begin{proof} 
For $\theta = (t-1)/t$, 
let $\mu = D^{1/\theta}_\# \nu_1$ and $ \tilde{\mu} = D^{1/\theta}_\# \tilde{\nu}_1$. By Theorem \ref{th:projextra} (point (1)), there exist $\theta$-strongly convex potentials $u$ and $\tilde{u}$ such that $\mathscr{P}^{cx}_{\preceq \nu_0}(\mu) = \nabla u^*_\# \nu_1$ and $\mathscr{P}^{cx}_{\preceq \nu_0}(\tilde{\mu}) = \nabla \tilde{u}^*_\# \tilde{\nu}_1$. 
Let $X_{\nu_1}, X_{\tilde{\nu}_1}\in L^2(\rho_0)$ be arbitrary maps such that $(X_{\nu_1})_\# \rho_0 = \nu_1$ and $(X_{\tilde{\nu}_1})_\# \rho_0 = \tilde{\nu}_1$. 
{Define
$\hbox{$X_\mu \coloneqq D^{1/\theta} \circ X_{\nu_1}$ and $X_{\tilde{\mu}} \coloneqq D^{1/\theta} \circ X_{\tilde{\nu}_1}$.}$
 Notice that  $T\coloneqq \nabla u^* \circ D^{\theta}$ and $\tilde{T} \coloneqq \nabla \tilde{u}^* \circ D^\theta$ are optimal transport maps  from $\mu$ to $\mathscr{P}^{cx}_{\preceq \nu_0}(\mu)$ and from $\tilde{\mu}$ to $\mathscr{P}^{cx}_{\preceq \nu_0}(\tilde{\mu})$, respectively (note in particular that $T$ and $\tilde{T}$ are the gradients of convex functions and they are 1-Lipschitz since $u$ and $\tilde{u}$ are} $\theta$-strongly convex).
Then, Lemma \ref{lem:lift} and equation \eqref{eq:stabproj} (the second inequality) imply that
\begin{align}\label{eq:extrastab1}
\|\nabla u^*\circ X_{\nu_1} - \nabla \tilde{u}^*\circ X_{\tilde{\nu}_1}\|_{L^2(\rho_0)}&  = \| T\circ X_\mu - \tilde{T}\circ X_{\tilde{\mu}} \|_{L^2(\rho_0)} \\\nonumber
& \le  \|  X_\mu -  X_{\tilde{\mu}} \|_{L^2(\rho_0)}
 = \frac{1}{\theta} \|X_{\nu_1} - X_{\tilde{\nu}_1}\|_{L^2(\rho_0)}\,.
\end{align}
Moreover, \[(tX_{\nu_1} + (1-t) \nabla u^*\circ X_{\nu_1},tX_{\tilde{\nu}_1} + (1-t)  \nabla \tilde{u}^*\circ X_{\tilde{\nu}_1})_\# \rho_0 \in \Pi(\mathscr{E}^t(\nu_0,\nu_1), \mathscr{E}^t(\nu_0,\tilde{\nu}_1))\,,\]
where $\Pi(\rho,\tilde{\rho})\in\mc{P}_2(\mathbb{R}^d\times \mathbb{R}^d)$ denotes the set of couplings with given marginals $\rho\in\mc{P}_2(\mathbb{R}^d)$ and $\tilde{\rho}\in\mc{P}_2(\mathbb{R}^d)$. Multiplying both sides of \eqref{eq:extrastab1} by $(t-1)$ (recalling $\theta = (t-1)/t$) and adding $t\|X_{\nu_1} - X_{\tilde{\nu}_1}\|_{L^2(\rho_0)}$ we get, by triangle inequality,
\[
{W_2(\mathscr{E}^t(\nu_0,\nu_1),\mathscr{E}^t(\nu_0,\tilde{\nu}_1))}\leq  2  t \|X_{\nu_1} - X_{\tilde{\nu}_1}\|_{L^2(\rho_0)}\,.
\]
Optimizing over $X_{\nu_1}$ and $ X_{\tilde{\nu}_1}$ we obtain the result. 
\end{proof}

\end{document}
